% Bagian: first-order-logic
% Bab: natural-deduction
% Bagian: derivations

\documentclass[../../../include/open-logic-section]{subfiles}

\begin{document}

\iftag{FOL}
      {\olfileid[id]{fol}{ntd}{der}}
      {\olfileid[id]{pl}{ntd}{der}}

\olsection{\usetoken{P}{derivation}}

\begin{explain}
Kita telah menjelaskan apa itu asumsi dan memberikan aturan-aturan
inferensi. !!^{derivation}s dalam deduksi alami dibangkitkan secara
induktif dari keduanya: setiap !!{derivation} merupakan sebuah asumsi
itu sendiri, atau terdiri atas satu, dua, atau tiga !!{derivation}s yang
diikuti oleh suatu inferensi yang benar.
\end{explain}

\begin{defn}[!!^{derivation}]
\Article{derivation} \emph{!!{derivation}} atas !!a{sentence}~$!A$ dari
asumsi-asumsi~$\Gamma$ adalah pohon berhingga yang terdiri atas
!!{sentence}s dan memenuhi syarat berikut:
\begin{enumerate}
\item !!^{sentence}s paling atas pada pohon itu berada dalam $\Gamma$
  atau !!{discharged} oleh suatu inferensi pada pohon tersebut.
\item !!^{sentence} paling bawah pada pohon itu adalah~$!A$.
\item Setiap !!{sentence} pada pohon itu selain kalimat~$!A$ di bagian
  paling bawah merupakan premis dari penerapan aturan inferensi yang
  benar, dan kesimpulan aturan tersebut berada tepat di bawah
  !!{sentence} itu pada pohon.
\end{enumerate}
Kita lalu mengatakan bahwa $!A$ adalah \emph{kesimpulan} dari
!!{derivation} tersebut dan $\Gamma$ adalah asumsi-asumsinya yang
!!{undischarged}.

Jika terdapat !!a{derivation} atas $!A$ dari~$\Gamma$, kita mengatakan
bahwa $!A$ \emph{!!{derivable}} dari~$\Gamma$, atau dalam simbol:
$\Gamma \Proves !A$. Jika terdapat !!a{derivation} atas~$!A$ yang
semua asumsinya !!{discharged}, kita menulis~$\Proves !A$.
\end{defn}

\begin{ex}
Setiap asumsi itu sendiri merupakan !!a{derivation}. Jadi, misalnya,
$!A$ sendiri merupakan !!a{derivation}, demikian pula $!B$ sendiri.
Kita dapat memperoleh !!a{derivation} baru dari keduanya dengan
menerapkan, misalnya, aturan $\Intro{\land}$,
\begin{prooftree}
\AxiomC{$!A$}
\AxiomC{$!B$}
\RightLabel{\Intro{\land}}
\BinaryInfC{$!A \land !B$}
\end{prooftree}
Aturan-aturan ini dimaksudkan untuk berlaku umum: kita dapat mengganti
$!A$ dan~$!B$ di dalamnya dengan sebarang !!{sentence}s, misalnya
dengan $!C$ dan~$!D$. Kesimpulannya lalu menjadi $!C \land !D$, sehingga
\begin{prooftree}
\AxiomC{$!C$}
\AxiomC{$!D$}
\RightLabel{\Intro{\land}}
\BinaryInfC{$!C \land !D$}
\end{prooftree}
merupakan !!a{derivation} yang benar. Tentu saja, kita juga dapat
menukar asumsi-asumsinya sehingga $!D$ memainkan peran~$!A$ dan $!C$
memainkan peran~$!B$. Jadi,
\begin{prooftree}
\AxiomC{$!D$}
\AxiomC{$!C$}
\RightLabel{\Intro{\land}}
\BinaryInfC{$!D \land !C$}
\end{prooftree}
juga merupakan !!a{derivation} yang benar.

Sekarang kita dapat menerapkan aturan lain, misalnya $\Intro{\lif}$,
yang memungkinkan kita menyimpulkan suatu implikasi dan !!{discharge}
asumsi apa pun yang identik dengan anteseden implikasi tersebut. Jadi,
kedua contoh berikut merupakan !!{derivation}s yang benar:
\begin{prooftree}
\AxiomC{$\Discharge{!C}{1}$}
\AxiomC{$!D$}
\RightLabel{\Intro{\land}}
\BinaryInfC{$!C \land !D$}
\DischargeRule{\Intro{\lif}}{1}
\UnaryInfC{$!C \lif (!C \land !D)$}
\DisplayProof\bottomAlignProof
\AxiomC{$!C$}
\AxiomC{$\Discharge{!D}{1}$}
\RightLabel{\Intro{\land}}
\BinaryInfC{$!C \land !D$}
\DischargeRule{\Intro{\lif}}{1}
\UnaryInfC{$!D \lif (!C \land !D)$}
\end{prooftree}
Keduanya masing-masing menunjukkan bahwa $!D \Proves !C \lif (!C
\land !D)$ dan $!C \Proves !D \lif (!C \land !D)$.

Ingatlah bahwa pelepasan asumsi merupakan izin, bukan kewajiban: kita
tidak harus melepaskan asumsi-asumsi tersebut. Secara khusus, kita dapat
menerapkan suatu aturan sekalipun asumsi yang dimaksud tidak terdapat
dalam !!{derivation}. Sebagai contoh, berikut ini diperbolehkan meskipun
tidak ada asumsi~$!A$ untuk !!{discharged}:
\begin{prooftree}
  \AxiomC{$!B$}
  \DischargeRule{\Intro{\lif}}{1}
  \UnaryInfC{$!A \lif !B$}
\end{prooftree}
\end{ex}

\end{document}
